% ---- pg-033 （书页 21） ----

验者完成的,不同选择会产生不同的结果.为了取得令人满意的表示常常需要一些重复步骤.

\begin{mvfigure}
\mvinc{pg-033}{0.212,0.121,0.819,0.445}
\phantomsection\label{fig:1.16}\mvcap{图 1.16\quad 前 5 家公用事业公司的星形}
\end{mvfigure}

切尔诺夫脸对核实以下两点很有帮助:(1)由题材知识和直觉知识提出的最初的分组;(2)由聚类算法产生的最终的分组.

\mvbold{例 1.12}\quad(切尔诺夫脸的实用数据)

根据\hyperref[tab:12.4]{表 12.4} 中的数据,22 家公用事业公司的观测值用切尔诺夫脸表示出来.我们得到以下的对应关系:

\begin{mvtable}
\begin{tabular}{ll}
\toprule
\textbf{变\quad 量} & \textbf{面部特征} \\
\midrule
$x_1$: 固定费用保障比率 & $\leftrightarrow$ 脸的半高度 \\
$x_2$: 资本回报率 & $\leftrightarrow$ 脸的宽度 \\
$x_3$: 每千瓦负载量的费用 & $\leftrightarrow$ 嘴中心的位置 \\
$x_4$: 每年负载因子 & $\leftrightarrow$ 眼睛的倾斜 \\
$x_5$: 从 1974 年起高峰期千瓦时需求的增长量 & $\leftrightarrow$ 眼睛的离心率(高/宽) \\
$x_6$: 销售额总量(每年使用千瓦时) & $\leftrightarrow$ 眼睛的半长 \\
$x_7$: 核的百分比 & $\leftrightarrow$ 嘴的曲率 \\
$x_8$: 总的燃料费用(美分/千瓦时) & $\leftrightarrow$ 鼻子长度 \\
\bottomrule
\end{tabular}
\end{mvtable}

\hyperref[fig:1.17]{图 1.17} 给出了切尔诺夫脸.我们已主观地将“相似的”脸分为 7 组.如果需要组的数目更少一些,我们可组合 5,6 和 7 组,甚至可将 2 和 3 组合并最终得到 4 或 5 个组.由于我们对与脸部容貌对应的变量的分配,公司的分组很大程度是依据地理位置.

完成构造切尔诺夫脸的工作必须有计算机的帮助.为了确定脸部特征的位置、尺寸和方向,在计算机程序中,通常把数据标准化作为过程的一部分.经过一些练习,我们可用切尔诺夫脸来传达相似性或不相似性,如下面的例子所说明的.
